
\lstset{
  basicstyle=\ttfamily\small,
  breaklines=true,
  columns=fullflexible,
  frame=single,
  keepspaces=true
}

\section{Lattice Mode for Emacs}

\subsection{Introduction}

This section describes the installation, configuration, and technical operation of a custom Emacs major mode designed for editing ELEGANT lattice files (typically \texttt{.lte} and \texttt{.lat}).  
The goal of this mode is to provide:
\begin{itemize}
\item Structured editing support for lattice definitions
\item Context-sensitive autocompletion of element types
\item Context-sensitive autocompletion of element parameters
\item Immediate access to parameter metadata via \texttt{=TAB}
\item Automatic activation for lattice file extensions
\end{itemize}

The implementation integrates tightly with Emacs’ native completion system and derives its semantic knowledge directly from ELEGANT’s source code.

\subsection{System Architecture}

The mode consists of two primary components:
\begin{enumerate}
\item \textbf{Major Mode Implementation}  
      \texttt{elegant-lattice-mode.el}

\item \textbf{Generated Data File}  
      \texttt{elegant-lattice-data.el}
\end{enumerate}

The data file is automatically generated from ELEGANT’s \texttt{track\_data.c} using a standalone Python script:

\begin{center}
\texttt{generate\_elegant\_lattice\_data\_el.py}
\end{center}

This separation ensures that:

\begin{itemize}
\item The editor remains synchronized with the ELEGANT source
\item Updates to ELEGANT require only regeneration of the data file
\item The Emacs mode remains stable and version-independent
\end{itemize}

\subsection{Data Extraction from \texttt{track\_data.c}}

The generator script parses several parallel arrays defined in \texttt{track\_data.c}:

\begin{itemize}
\item \texttt{entity\_name[]} — element type identifiers
\item \texttt{entity\_text[]} — short textual descriptions
\item \texttt{entity\_description[]} — parameter array linkage
\item \texttt{<type>\_param[]} — parameter definitions
\end{itemize}

For each element type, the script performs:

\begin{enumerate}
\item Extraction of the canonical element name
\item Retrieval of the descriptive text
\item Identification of the associated parameter array
\item Parsing of each parameter entry to obtain:
      \begin{itemize}
      \item Parameter name
      \item Units
      \item Parameter type (if present)
      \item Description string
      \end{itemize}
\end{enumerate}

The script then emits four Emacs Lisp variables:

\begin{itemize}
\item \texttt{elegant-lattice-element-types}
\item \texttt{elegant-lattice-element-docs}
\item \texttt{elegant-lattice-element-params}
\item \texttt{elegant-lattice-param-docs}
\end{itemize}

These structures form the semantic backend for both completion and documentation display.

\section{Installation Procedure}

\subsection{Directory Layout}

It is recommended to install the mode in a dedicated directory:

\begin{lstlisting}
~/.emacs.d/elegant-lattice/
\end{lstlisting}

Place the following files in this directory:

\begin{lstlisting}
elegant-lattice-mode.el
elegant-lattice-data.el
\end{lstlisting}

\subsubsection{Generating Completion Data}

Generate the data file using:

\begin{lstlisting}
python3 generate_elegant_lattice_data_el.py /path/to/track_data.c \
  -o ~/.emacs.d/elegant-lattice/elegant-lattice-data.el
\end{lstlisting}

This step should be repeated whenever ELEGANT is updated.

\subsubsection{Emacs Configuration}

Add the directory to your \texttt{load-path} and require the mode:

\begin{lstlisting}
(add-to-list 'load-path (expand-file-name "~/.emacs.d/elegant-lattice"))
(require 'elegant-lattice-mode)
\end{lstlisting}

Ensure automatic activation:

\begin{lstlisting}
(add-to-list 'auto-mode-alist '("\\.lte\\'" . elegant-lattice-mode))
(add-to-list 'auto-mode-alist '("\\.lat\\'" . elegant-lattice-mode))
\end{lstlisting}

Restart Emacs after editing your initialization file.

\subsection{Completion Behavior}

The mode integrates with Emacs’ native \texttt{completion-at-point} (CAPF) mechanism.

\subsubsection{Element Type Completion}

Example:

\begin{lstlisting}
Q1: qu
\end{lstlisting}

Invoke completion with:

\begin{lstlisting}
M-TAB
\end{lstlisting}

\subsubsection{Parameter Completion}

Example:

\begin{lstlisting}
Q1: quad, k
\end{lstlisting}

Invoke:

\begin{lstlisting}
M-TAB
\end{lstlisting}

Completion is case-insensitive but preserves user typing style.

\subsection{Parameter Documentation Using \texttt{=TAB}}

The mode provides contextual documentation lookup for parameter assignments.

Example:

\begin{lstlisting}
Q1: quad, k1=
\end{lstlisting}

Placing the cursor immediately after the equals sign and pressing:

\begin{lstlisting}
TAB
\end{lstlisting}

produces echo-area output such as:

\begin{lstlisting}
QUAD.K1  units=1/m^2  type=double  normalized quadrupole strength
\end{lstlisting}

This mechanism:

\begin{itemize}
\item Extracts the element type from the current line
\item Identifies the parameter preceding the equals sign
\item Looks up metadata in \texttt{elegant-lattice-param-docs}
\item Displays units, parameter type, and description
\end{itemize}

If metadata is unavailable, a minimal fallback message is displayed.

This feature enables rapid inspection of parameter meaning without leaving the editing buffer.

\subsection{Optional Popup Completion Interfaces}

For graphical Emacs users, popup completion frontends may be used.

\subsubsection*{Corfu}

\begin{lstlisting}
(require 'corfu)
(global-corfu-mode 1)
(setq corfu-auto t)
(setq corfu-cycle t)
\end{lstlisting}

\subsubsection*{Company}

\begin{lstlisting}
(require 'company)
(global-company-mode 1)
(setq company-idle-delay 0.0)
(setq company-minimum-prefix-length 1)
\end{lstlisting}

Both integrate seamlessly with the mode’s CAPF backend.

\subsection{Key Binding Summary}

\begin{center}
\begin{tabular}{ll}
\toprule
\textbf{Key} & \textbf{Function} \\
\midrule
\texttt{M-TAB} & Complete element type or parameter \\
\texttt{TAB} after \texttt{=} & Show parameter metadata \\
\texttt{TAB} elsewhere & Indent line \\
\bottomrule
\end{tabular}
\end{center}

\subsection{Conclusion}

This Emacs lattice mode provides a structured, source-synchronized editing environment for ELEGANT lattice files. By
deriving semantic information directly from \texttt{track\_data.c}, it ensures long-term maintainability and version
consistency while significantly improving editing efficiency.

